% section: 共通する変形

\section{共通する変形}

個別のパターンは,ばらばらの技ではない.多くの場合,使っている
操作は「規則を保つ変数を探す」という一つの考え方にまとめられる.
ここでは,パターン集を読む前に,頻繁に現れる変形を導出しておく.

\subsection{不動点を引く}

\[
  a_{n+1}=pa_n+q
\]
で,定数項 $q$ が邪魔をしているとする.数列が一定になる値 $\alpha$ を
探すと,
\[
  \alpha=p\alpha+q
\]
である.ここで $b_n=a_n-\alpha$ と置けば,
\begin{align*}
  b_{n+1}
    &=a_{n+1}-\alpha\\
    &=pa_n+q-(p\alpha+q)\\
    &=p(a_n-\alpha)=pb_n.
\end{align*}
不動点を引くとは,定数を消して等比型へ移すことである.

\subsection{補助数列で非斉次項を消す}

\[
  a_{n+1}=pa_n+f(n)
\]
では, $f(n)$ と同じ種類の項を含む数列 $g(n)$ を探し,
\[
  g(n+1)=pg(n)+f(n)
\]
を満たすようにする. $b_n=a_n-g(n)$ と置けば,
\[
  b_{n+1}=pb_n
\]
となる.この $g(n)$ は,元の漸化式の特別な一つの解を先に探して
いると考えられる.残りが等比型になるから,「特解を引く」という操作が
有効なのである.

\subsection{差と比を,それぞれ和と積へ戻す}

差分
\[
  a_{n+1}-a_n=f(n)
\]
を $n=1$ から $n-1$ まで足すと,
\[
  a_n-a_1=\sum_{k=1}^{n-1}f(k)
\]
を得る.これは離散的な積分である.同様に,
\[
  \frac{a_{n+1}}{a_n}=g(n)
\]
を掛け合わせると,
\[
  a_n=a_1\prod_{k=1}^{n-1}g(k)
\]
となる.差には和,比には積が対応するという対応を覚えておくと,
基本パターンの見分けが速くなる.

\subsection{逆数と不動点で分数をほどく}

\[
  a_{n+1}=\frac{pa_n}{qa_n+r}
\]
では,逆数 $b_n=1/a_n$ を取ると,
\[
  b_{n+1}=\frac{q}{p}+\frac{r}{p}b_n
\]
となる.分数が消え,一次漸化式に変わる.

より一般に,
\[
  a_{n+1}=\frac{pa_n+q}{ra_n+s}
\]
では,不動点 $\alpha,\beta$ を使って
\[
  b_n=\frac{a_n-\alpha}{a_n-\beta}
\]
と置くと,うまくいく場合がある.この比が等比型になるのは,
二つの不動点からの距離の比が,一次分数変換によって一定倍率で変化
するからである.詳しい意味は付録の分数一次変換で扱う.

\subsection{特性方程式と行列}

隣接する過去の項が線形に現れるとき,
\[
  a_{n+2}=pa_{n+1}+qa_n
\]
に対して $a_n=\lambda^n$ を試す.すると,
\[
  \lambda^2-p\lambda-q=0
\]
が得られる.指数関数型の基本解を探しているので,この方程式を
特性方程式と呼ぶ.

同じ規則は行列で
\[
  \begin{pmatrix}a_{n+2}\\a_{n+1}\end{pmatrix}
  =
  \begin{pmatrix}p&q\\1&0\end{pmatrix}
  \begin{pmatrix}a_{n+1}\\a_n\end{pmatrix}
\]
とも書ける.特性方程式は,この行列が持つ固有値を調べているのである.

\subsection{恒等式で非線形性を翻訳する}

\[
  a_{n+1}=2a_n^2-1
\]
は, $a_n=\cos\theta_n$ と置けば $\theta_{n+1}=2\theta_n$ になる.
\[
  a_{n+1}=a_n^2+2
\]
は, $a_n=\alpha^{2^{n-1}}-\alpha^{-2^{n-1}}$ と置けば,
同じ形が保たれる.このような置換は,非線形な式を無理に一次化する
のではなく,非線形性が自然に現れる別の変数へ移している.

第2章第7節の特殊関数型では,この操作を解法として使った.付録の
チェビシェフ多項式,特殊関数,力学系では,その操作がなぜ他の問題にも
現れるのかをさらに広い視点から調べる.
